\section{Convex support on endpoint caps}
\label{sec:v75-convex-support}
Put $R_-=9/20$, $R_+=47/100$ and $\tau_-=3/50$.
A word $w=(C_0,\ldots,C_m)$ records the actual successive disk
centers, with $C_0=0$. A regular word has positive flights, positive
incidences, and the original first-hit exclusions. Its physical
endpoint image need not be connected or convex.

For a regular trajectory write $Q_i=C_i+Rn_i$, let $v_i$ be its
outgoing unit velocity, and set
\[
 w_0=v_0,\qquad w_m=-v_{m-1},\qquad
 n_0\cdot w_0=c_0>0,\quad n_m\cdot w_m=c_m>0.
\]
Here $v_{m-1}$ is the velocity on the last flight, before its terminal
reflection. Thus the endpoint gradients of the polygonal length are
$-w_0$ and $-w_m$.

\begin{lemma}[A fixed cover of regular endpoint data]
\label{lem:v75-cap-cover}
Let $e_j=(\cos(2\pi j/16),\sin(2\pi j/16))$, $0\le j<16$.
For every pair of unit vectors $n,w$ with $n\cdot w>0$, some $j$
satisfies $n\cdot e_j>1/2$ and $w\cdot e_j>1/2$.
On that cap the coordinate
\begin{equation}\label{eq:v75-cap-coordinate}
 Q_j(s)=C+s e_j^\perp+\sqrt{R^2-s^2}\,e_j,
 \qquad |s|<\sqrt3 R/2,
\end{equation}
is single-valued, with $|Q_j'|\le2$ and $|Q_j''|\le8/R_-$.
Choosing the first admissible $j$ at each endpoint gives a Borel
partition with at most $16^2=256$ parts, independent of $m$ and $R$.
\end{lemma}
\begin{proof}
The shorter angular distance between $n$ and $w$ is less than
$\pi/2$. Its midpoint is within $\pi/4$ of both. A closest grid
direction is within $\pi/16$ of the midpoint; hence both distances
are less than $5\pi/16<\pi/3$. The two scalar products exceed $1/2$.
Formula \eqref{eq:v75-cap-coordinate} parameterizes the stated
semicircle cap. Its derivatives are
\[
 Q_j'=e_j^\perp-\frac{s}{\sqrt{R^2-s^2}}e_j,
 \qquad Q_j''=-\frac{R^2}{(R^2-s^2)^{3/2}}e_j.
\]
These give the bounds and the measurable first-admissible partition.
No lower bound on the physical incidence $n\cdot w$ is used.
\end{proof}

\subsection{Supporting actions, not a continuation of the orbit}
For the prescribed word define the auxiliary endpoint value
\begin{equation}\label{eq:v75-convex-value}
 V_w(q_0,q_m)=\min_{q_i\in\overline{D(C_i,R)},\ 0<i<m}
                         \sum_{i=0}^{m-1}|q_{i+1}-q_i|.
\end{equation}
For $m=1$ this is just $|q_m-q_0|$. The minimization is a comparison
in Euclidean space. Its other minimizers are not declared to be
billiard trajectories; no measure or label is assigned to them.

\begin{lemma}[A global supporting inequality]
\label{lem:v75-support}
At a regular physical trajectory $x$ of word $w$, with roof $L(x)$,
\begin{align}
 V_w(q_0,q_m)&\ge \ell_x(q_0,q_m),\notag\\
 \ell_x(q_0,q_m)&=L(x)-w_0(x)\cdot(q_0-Q_0(x))
                         -w_m(x)\cdot(q_m-Q_m(x)).
                                      \label{eq:v75-support-affine}
\end{align}
Equality holds at its own endpoints. A regular physical trajectory
of the word is uniquely determined by its endpoints.
\end{lemma}
\begin{proof}
The polygonal length is jointly convex in all contact positions.
At an internal reflected contact its gradient is $-2c_i n_i$.
For $q_i'$ in the corresponding closed disk,
$-2c_i n_i\cdot(q_i'-Q_i)\ge0$, strictly if $q_i'\ne Q_i$,
because $c_i>0$ and the disk is strictly convex. The convex gradient
inequality, followed by minimization over the internal disks, gives
\eqref{eq:v75-support-affine}. With the endpoints unchanged it also
shows that the physical internal tuple is the unique minimizer.
Two regular trajectories of the same word and endpoints must therefore
coincide. The one-flight case follows from its straight segment.
The minimum in \eqref{eq:v75-convex-value} exists by compactness.
It is convex in the Cartesian endpoints: interpolate both endpoint
pairs and their minimizing internal tuples, and apply convexity of
the polygonal length. The supporting inequality, rather than any
physical interpretation of interpolated tuples, is what will be used.
\end{proof}

Fix a pair $a=(j_0,j_m)$ of admissible caps and let $\mathcal U_{w,a}$
be the full regular physical endpoint image satisfying both strict
cap tests. It is considered before the first-admissible partition
and before restriction to any return label or source weight. In the
open rectangle $Q_R=(-\sqrt3R/2,\sqrt3R/2)^2$ define
\begin{equation}\label{eq:v75-support-envelope}
 \Phi_{w,a}(s_0,s_m)=
 \sup_{x\in\mathcal U_{w,a}}
             \ell_x\bigl(Q_{j_0}(s_0),Q_{j_m}(s_m)\bigr).
\end{equation}
An empty $\mathcal U_{w,a}$ is omitted.

\begin{theorem}[Uniform convex support for every physical word]
\label{thm:v75-convex-envelope}
The function \eqref{eq:v75-support-envelope} is finite and
$\lambda$-strongly convex on $Q_R$, where
\begin{equation}\label{eq:v75-uniform-curvature}
 \lambda=\frac1{2R_+}=\frac{50}{47},\qquad
 \Lambda=\frac8{\tau_-}+\frac8{R_-}=\frac{1360}{9}.
\end{equation}
It equals the actual reduced roof $F_w$ on $\mathcal U_{w,a}$.
On every regular open piece of that image,
\begin{equation}\label{eq:v75-hessian-sandwich}
                  \lambda I_2\le D^2F_w\le\Lambda I_2.
\end{equation}
Both constants are independent of the word length, every internal
incidence, and every positive physical clearance. No uniform size
of a physical neighborhood and no bound on higher derivatives is
asserted.
\end{theorem}
\begin{proof}
Each supporting function in \eqref{eq:v75-support-envelope} has
zero mixed second derivative in these coordinates. At an endpoint
its second derivative is
\[
  (w_i(x)\cdot e_{j_i})\frac{R^2}{(R^2-s_i^2)^{3/2}}
                                  \ge\frac1{2R_+}.
\]
Subtracting $\lambda|s|^2/2$ therefore makes every support convex;
their supremum is convex. For fixed $m$ all support values are
bounded, since $L(x)\le m\tau_+$ and the endpoint displacements
are at most $2R$ each. The envelope is finite.
By Lemma~\ref{lem:v75-support} every support is at most
$V_w(Q_{j_0}(s_0),Q_{j_m}(s_m))$. At a physical point its own
support attains the roof, which equals that value. Thus the envelope
agrees with the roof throughout the open physical cap image.
Endpoint injectivity follows from the same lemma. Local smoothness
of the reduced action follows from the regular reflection equations
and the positive definite internal contact Hessian; this is also
the endpoint construction in the inherited curvature lemma.
Strong convexity of the envelope proves the lower inequality.

For the upper inequality one can keep all internal contacts fixed
while the two endpoints vary. The resulting trial length touches
the reduced action from above. Equivalently, the reduced Hessian is
the Schur complement of the internal contact Hessian. For an endpoint
variation $h=(h_0,h_m)$ the transverse second variation of a segment
is its squared transverse displacement divided by its flight length.
At most two endpoint displacements occur. Using $|Q_j'|\le2$ and
$\ell_i\ge\tau_-$ bounds their total by $8|h|^2/\tau_-$,
including $m=1$, when they are on the same segment. The endpoint
curvature terms are at most $8|h|^2/R_-$. Minimizing over internal
variations cannot increase this trial quadratic form. This proves
the upper inequality. The argument is local at a regular trajectory;
it makes no continuation across a grazing or first-hit boundary.
\end{proof}

The use of supporting actions is essential. Positivity of a Hessian
on a disconnected physical image alone would not bound the total
turning of all its level components. The single convex envelope on
the entire endpoint rectangle supplies that missing global comparison.
